\section{Speed-shell cancellation and the effective defect}
\label{ss:section}

\paragraph{Status and scope.}
This is a full component derivation, author-checked and awaiting independent
mathematical audit. It does not establish an arbitrary-data critical time
bound. The snapshot assumptions are the accepted quotient identities for
solenoidal $u\in H^m(\mathbb R^3)^3$, $m\ge4$: with
$w=u+q$, $\rho=|w|$, and $\sigma=-\operatorname{div}w$,
\[
 w\in L^3\cap L^6,\quad \nabla w\in L^2,\quad
 \operatorname{div}(\rho w)=0,\quad w\cdot\nabla\rho=\rho\sigma.
\]
The last identity includes the zero set, on which $\sigma=0$ almost
everywhere. No $L^2$ hypothesis on $w$, spatial moment, smoothness of the
representative, or differentiation of its unit direction is assumed.
The time-dependent statements concern only the original unforced equation
on $\mathbb R^3$, viscosity $\nu>0$, and the selected Schwartz-data branch.
We use $V=\rho^{1/2}w$, $Q=\|w\|_3^3/3$, $D=D_{\mathcal Q}$,
$a_0=9C_S^2/8$, and $B^0=(R_iR_j\sigma)+\sigma I/3$ as in
Section~\ref{sd:section}.

\subsection{The defect has zero average on every speed shell}

\begin{theorem}[Renormalized shell cancellation]\label{ss:shell}
For every $k>0$ the integral below is absolutely convergent and
\begin{equation}\label{ss:superlevel}
 \int_{\{\rho>k\}}\sigma\,dx=0.
\end{equation}
In particular $\int_{\{a<\rho\le b\}}\sigma=0$ for $0<a<b<\infty$.
Let $\mathcal H_\rho$ be the closed subspace of real $L^2(\mathbb R^3)$
generated by these shell indicators with positive rational endpoints.
Then
\begin{equation}\label{ss:orthogonal}
 \sigma\perp\mathcal H_\rho.
\end{equation}
Every Borel $g:(0,\infty)\to\mathbb R$ for which $g(\rho)$, extended
by zero on $\{\rho=0\}$, belongs to $L^2$ gives an element of
$\mathcal H_\rho$ and hence $\int\sigma g(\rho)=0$.
\end{theorem}
\begin{proof}
Use the locally Sobolev vector field
\[
 F_k=(\rho/k-1)_+w.
\]
The chain rule, or a smooth truncation followed by the local Sobolev
chain rule, gives
\[
 \operatorname{div}F_k
 =-(\rho/k-1)_+\sigma+
   k^{-1}\mathbf1_{\{\rho>k\}}w\cdot\nabla\rho
 =\mathbf1_{\{\rho>k\}}\sigma.
\]
There is no surface measure at $\rho=k$: the positive-part map is
Lipschitz, and $\nabla\rho=0$ almost everywhere on that level set.
The derivatives are locally integrable, since their magnitudes are
bounded by $C(1+\rho/k)|\nabla w|$ on bounded sets.
Also $|F_k|\le\rho^2/k$, so $F_k\in L^{3/2}$, whereas
\[
 |\{\rho>k\}|\le k^{-3}\|w\|_3^3,\qquad
 \|\mathbf1_{\{\rho>k\}}\sigma\|_1
 \le |\{\rho>k\}|^{1/2}\|\sigma\|_2<\infty.
\]
For a cutoff $\chi_R$ equal to one on the ball of radius $R$, zero
outside the ball of radius $2R$, and with $|\nabla\chi_R|\le C/R$,
the norm $\|\nabla\chi_R\|_3$ is uniformly bounded. Thus
\[
 \left|\int F_k\cdot\nabla\chi_R\right|
 \le C\|F_k\|_{L^{3/2}(\{|x|>R\})}\longrightarrow0.
\]
Dominated convergence on the divergence proves \eqref{ss:superlevel}
for every $k>0$, not just for almost every threshold. Taking differences
and then $L^2$ closure proves \eqref{ss:orthogonal}.

For the final assertion, first restrict the speed to $[a,b]\subset
(0,\infty)$. The pushforward of Lebesgue measure restricted to
$\{a\le\rho\le b\}$ is a finite Borel measure on that interval.
Rational interval step functions are dense in its $L^2$ space.
Approximation and then exhaustion $a\downarrow0$, $b\uparrow\infty$
prove the assertion. This argument defines the subspace without taking
conditional expectation on an infinite-mass probability space. Values on
$\{\rho=0\}$ are fixed to zero throughout.
\end{proof}

\subsection{Only variation not determined by speed contributes}

Set
\[
 F^0=V\otimes V-\tfrac13|V|^2I,\qquad
 \chi=\sum_{i,j}R_iR_jF^0_{ij}
      =\sum_{i,j}R_iR_j(V_iV_j)+\tfrac13\rho^3.
\]
Here $\chi$ is a scalar $L^2$ field, not the velocity pressure. Let
$P_\rho$ be the Hilbert-space orthogonal projection onto $\mathcal H_\rho$
and define
\begin{equation}\label{ss:delta}
 N_\rho=\|(I-P_\rho)\chi\|_2,
 \qquad
 \delta_\rho=\frac{3N_\rho}{2\|\rho^3\|_2}
 \quad(w\ne0),
\end{equation}
with $\delta_\rho=0$ for $w=0$.

\begin{proposition}[Effective-defect estimate]\label{ss:residual}
One has $0\le\delta_\rho\le1$ and the exact identity and bound
\begin{align}
 K_u&=\langle\sigma,(I-P_\rho)\chi\rangle,\label{ss:pairing}\\
 |K_u|&\le\frac23\delta_\rho\|\sigma\|_2\|w\|_6^3.
       \label{ss:effective}
\end{align}
For any finite disjoint family $E_j=\{a_j<\rho\le b_j\}$,
$0<a_j<b_j<\infty$, zero-volume terms omitted,
\begin{equation}\label{ss:finite-shell}
 N_\rho^2\le\|\chi\|_2^2-
     \sum_j\frac{\bigl(\int_{E_j}\chi\,dx\bigr)^2}{|E_j|}.
\end{equation}
The right side decreases to $N_\rho^2$ for nested dyadic speed
partitions exhausting $(0,\infty)$. Thus finite shell averages give
certified upper bounds for the analytic residual when their integrals
are known exactly; numerical quadrature alone is not such a certificate.
\end{proposition}
\begin{proof}
The signed identity of Lemma~\ref{sd:cancellation}, the $L^2$
self-adjointness of $R_iR_j$, and the trace-free nature of $B^0$ give
\[
 K_u=\langle B^0,F^0\rangle
     =\left\langle\sigma,\sum_{i,j}R_iR_jF^0_{ij}\right\rangle
     =\langle\sigma,\chi\rangle.
\]
All pairings are $L^2$ pairings. At a nonzero Fourier frequency the
restriction of $F\mapsto\sum R_iR_jF_{ij}$ to trace-free symmetric
matrices has norm $\sqrt{2/3}$: its symbol is contraction with the
trace-free part of $-\xi\otimes\xi/|\xi|^2$.
Pointwise $|F^0|_F=\sqrt{2/3}\rho^3$. Therefore
\[
 \|\chi\|_2\le\sqrt{2/3}\|F^0\|_2
              =\tfrac23\|\rho^3\|_2.
\]
Orthogonal projection is contractive, and Theorem~\ref{ss:shell}
annihilates $P_\rho\chi$. These facts prove all assertions up to
\eqref{ss:finite-shell}, including the case $\sigma=0$ without any
division by its norm.

For the finite family, the functions
$|E_j|^{-1/2}\mathbf1_{E_j}$ form an orthonormal set in
$\mathcal H_\rho$. The squared norm removed by their projection is
exactly the displayed sum. Pythagoras proves the inequality.
For a concrete nested family use mesh $2^{-m}$ on
$(2^{-m},2^m]$, with $m=1,2,\ldots$. The associated spaces are nested,
and their union is dense in $\mathcal H_\rho$ by the finite-measure
approximation in the preceding proof. Their projection norms therefore
increase to $\|P_\rho\chi\|_2$, proving the limit.
\end{proof}

\begin{proposition}[Measurability and a depleted fourth-power bound]
\label{ss:clock}
On the classical branch, $\delta_\rho(t)$ is Borel and
\begin{equation}\label{ss:depleted-diff}
 Q'+\frac\nu2D\le
  \frac{a_0^3}{2\nu^3}\delta_\rho^4\|\sigma\|_2^4Q.
\end{equation}
With
$L_\rho(t)=a_0^3(2\nu^3)^{-1}
\int_0^t\delta_\rho(s)^4\|\sigma(s)\|_2^4\,ds$,
\begin{equation}\label{ss:depleted-budget}
 Q(t)+\frac\nu2\int_0^tD(s)\,ds\le Q(0)e^{L_\rho(t)}.
\end{equation}
These are unconditional inequalities on each compact classical interval;
no endpoint-uniform upper bound for $L_\rho$ is asserted.
\end{proposition}
\begin{proof}
The accepted continuity $w(t)\in C_tL^3_x$ and local uniform $L^6$ bound
imply that $F^0(t)$ is continuous in $L^1$ and bounded in $L^2$ on each
compact time interval. Testing first with smooth compactly supported
functions, then using $L^2$ density, proves weak $L^2$ continuity.
Consequently $\chi(t)$ is weakly $L^2$ continuous. The same reasoning
applies to $\rho(t)^3$.

Choose the countable set $\mathcal C$ of continuous piecewise-linear
scalar functions, zero outside a compact subinterval of $(0,\infty)$,
with rational breakpoints and rational values, including the zero
function. For every fixed $\rho$, its compositions are dense in
$\mathcal H_\rho$; this follows from continuous-function approximation
for the finite pushforward measures used above. For fixed $g\in\mathcal C$,
$g(\rho(t))$ is strongly continuous in $L^2$: the union of the supports
at two times has uniformly bounded measure by the $L^3$ bound and the
positive lower endpoint of the support of $g$, and H\"older applied to
$|g(\rho(t))-g(\rho(s))|\le\operatorname{Lip}(g)|w(t)-w(s)|$
finishes the assertion. Thus
\[
 N_\rho(t)=\inf_{g\in\mathcal C}\|\chi(t)-g(\rho(t))\|_2
\]
is a countable infimum of lower semicontinuous functions, hence Borel.
The denominator in \eqref{ss:delta} is Borel by weak continuity. This
proves the measurability assertion, including the zero-denominator set.
The weak $L^2$ continuity of $\sigma$ proved in Section~\ref{sd:section}
shows that the whole integrand defining $L_\rho$ is Borel and locally
bounded, since $\|\sigma\|_2^2\le Y/4$ and $\delta_\rho\le1$.

Interpolation and the accepted $\|w\|_9^3\le a_0D$ give
\[
 |K_u|\le\tfrac23\delta_\rho\|\sigma\|_2
                    (3Q)^{1/4}(a_0D)^{3/4}.
\]
Use $\alpha x^{3/4}\le\varepsilon x+
27\alpha^4/(256\varepsilon^3)$ with $\varepsilon=\nu/2$ in
$Q'+\nu D=K_u$. The coefficient is exactly
$(27/32)(2/3)^4\,3a_0^3=a_0^3/2$.
The integrating factor $e^{-L_\rho}$, with its nonnegative dissipation
term retained, gives \eqref{ss:depleted-budget}.
\end{proof}

\subsection{Null scalar corrections to the signed form}

For the countable class $\mathcal C$ from the preceding proof set
\begin{equation}\label{ss:gauge-definition}
 \beta_\nu(t)=\inf_{h\in\mathcal C}
 b_\nu\bigl(B^0(t)+h(\rho(t))\sigma(t)I\bigr),
\end{equation}
where $b_\nu$ is exactly the nonnegative Rayleigh rate in
\eqref{sd:form-definition}, with Dirichlet coefficient $4\nu/9$.

\begin{proposition}[Speed-dependent null corrections]\label{ss:gauge}
The rate $\beta_\nu$ is Borel, locally bounded, and satisfies
\[
 0\le\beta_\nu\le b_\nu(B^0),\qquad
 Q'+\tfrac\nu2D\le3\beta_\nu Q.
\]
No optimizing function $h$ or measurable selection is assumed to exist.
\end{proposition}
\begin{proof}
For every bounded Borel $h$, the field $h(\rho)\rho^3$ is in
$\mathcal H_\rho$ by Theorem~\ref{ss:shell}, since $\rho^3\in L^2$.
Hence the exact signed work is unchanged:
\[
 \int V^{\mathsf T}\bigl(B^0+h(\rho)\sigma I\bigr)V=K_u.
\]
For $h\in\mathcal C$ this matrix lies in $L^2$. Its Rayleigh rate is
finite: Sobolev interpolation gives, for normalized tests,
$\int\psi^{\mathsf T}M\psi\le
C_S^{3/2}\|M\|_2\|\nabla\psi\|_2^{3/2}$, to which the same Young
maximization applies. The form inequality extends to $V\in H^1$ by
density. Since $D\ge(8/9)\|\nabla V\|_2^2$, it gives
$Q'+\nu D/2\le3b_\nu(B^0+h(\rho)\sigma I)Q$ for each $h$.
Taking the infimum proves the differential assertion. The zero profile
belongs to $\mathcal C$, so $0\le\beta_\nu\le b_\nu(B^0)$.

For each fixed smooth compactly supported vector test $\psi$,
$h(\rho(t))|\psi|^2$ is strongly continuous in $L^2$, whereas
$\sigma(t)$ is weakly continuous in $L^2$. Their pairing is continuous.
The $B^0$ term has the same property. As in Theorem~\ref{sd:form-clock},
a fixed countable dense set of vector tests gives a Borel Rayleigh
supremum. Taking the countable infimum over $h$ proves that $\beta_\nu$
is Borel. Its upper bound by the old rate gives local boundedness.
\end{proof}

\subsection{The exact remaining time estimate}

\begin{corollary}[Combined clock and the original missing bound]
\label{ss:consumer}
Define
\[
 c_\nu(t)=\min\left\{3\beta_\nu(t),
    \frac{a_0^3}{2\nu^3}\delta_\rho(t)^4\|\sigma(t)\|_2^4\right\},
 \qquad L_c(t)=\int_0^t c_\nu(s)\,ds.
\]
Then, for every $t<T_*$,
\begin{align}
 Q(t)+\frac\nu2\int_0^tD&\le Q_0e^{L_c(t)},\label{ss:combined}\\
 \int_0^t\|\sigma(s)\|_2^4\,ds
 &\le\frac{E_0Y_0}{32\nu}
       \exp\!\bigl(A_*e^{L_c(t)}\bigr),\label{ss:missing}\\
 A_*&=\frac{8C_S^3C_9^3Q_0}{3\nu^3}.\nonumber
\end{align}
Here the second implication uses the independently review-pending
Proposition~\ref{de:quotient-clock}. An input-only bound on $L_c$,
uniform below $\min\{H,T_*\}$ for every finite $H$, would therefore
supply the requested fourth-power estimate and global continuation.
\end{corollary}
\begin{proof}
Both preceding differential bounds hold almost everywhere and have the
same left side, so their minimum also bounds that side. Integrating gives
\eqref{ss:combined}. The quotient clock yields
\[
 \sup_{s\le t}Y(s)\le Y_0
  \exp\left(\frac{4C_S^3C_9^3}{3\nu^2}\int_0^tD\right)
 \le Y_0\exp(A_*e^{L_c(t)}).
\]
Combine $\|\sigma\|_2^4\le Y^2/16$ and
$\int_0^tY\le E_0/(2\nu)$ to prove \eqref{ss:missing}.
The local $H^1$ blow-up alternative supplies the continuation implication.
\end{proof}

\paragraph{Attempted closure and the first unsupported step.}
The exact shell identity eliminates a whole subspace of the signed pairing,
not merely one trace term. Nevertheless $\delta_\rho\le1$ only recovers
$L_c\le a_0^3(2\nu^3)^{-1}\int\|\sigma\|_2^4$, whose right side is
the original unknown. The energy and moment budgets established earlier
bound lower time exponents and do not, by that comparison, bound $L_c$.
No estimate on the time evolution of $P_\rho$, $\delta_\rho$ or an
optimizing form correction has been obtained here. Differentiating their
variational definitions without such a theorem would be an unsupported
step; none of the proofs above does so.

The depleted coefficient has retained, not improved, critical scaling.
For $u_\lambda(x)=\lambda u(\lambda x)$, quotient covariance gives
$\sigma_\lambda(x)=\lambda^2\sigma(\lambda x)$ and
$\chi_\lambda(x)=\lambda^3\chi(\lambda x)$.
Composition with the dilation identifies the closed speed subspaces,
so $\delta_{\rho_\lambda}=\delta_\rho$ and
$\delta_{\rho_\lambda}^4\|\sigma_\lambda\|_2^4
=\lambda^2\delta_\rho^4\|\sigma\|_2^4$.
With the Navier--Stokes time scaling its integral is invariant.
No strict separation of finiteness conditions on actual maximal branches,
new small-data theorem, or universal depletion bound is asserted.

The related one-sided strain and sum-space literature, including
\cite{Miller2021Sum}, is background rather than an input supplying the
missing bound: its velocity-strain hypotheses are not estimates for the
present speed residual or correction Hessian. These component proofs and
all previously pending components still require independent review.
The arbitrary-data temporal producer and the terminal claim remain open.
